\newpage
\section{CONCLUSIONES Y RECOMENDACIONES}
A lo largo del desarrollo del informe se pusieron a prueba una serie de conceptos y prácticas importantes para el manejo básico del ESP8266, para la construcción del sistemas de suministro de comida y agua para animales, utilizando IoT. A continuación algunas conclusiones y recomendaciones recolectadas a lo largo de trabajo realizado: 

\begin{itemize}
    \item A lo largo del proyecto, se realizaron ajustes importantes en el diseño y selección de componentes para cumplir con las prioridades definidas. La adaptación a las limitaciones técnicas, como el cambio de motores y la búsqueda de alternativas, permitió avanzar en el desarrollo del dispensador de comida y agua, logrando el objetivo principal del proyecto.

    \item El proyecto implicó enfrentar varios problemas técnicos, como la falta de potencia del ESP8266 y dificultades con Arduino Cloud. Se aplicaron estrategias efectivas, como la transición a Blynk, para superar estos obstáculos y garantizar el funcionamiento del sistema de dispensación de comida y agua, demostrando flexibilidad y capacidad de resolución.
    
    \item Se logró implementar funcionalidades básicas de IoT mediante la integración de Blynk, lo que permitió el monitoreo remoto del estado de los recipientes y el control del sistema desde un dispositivo móvil. Esta implementación cumplió parcialmente con el objetivo de automatización del suministro y facilitó el control a distancia del sistema.

    \item La construcción del hardware, que incluye el dispensador de comida, el sistema de bombas para agua y los sensores ultrasónicos, fue completada exitosamente. Este aspecto del proyecto proporcionó una base sólida para el funcionamiento básico del sistema, permitiendo que el dispensador cumpliera con sus funciones principales de manera efectiva.

    \item Para mejorar el rendimiento del servomotor y manejar mayores cargas, se recomienda replicar el circuito de relé. Este ajuste permitirá al servomotor operar con mayor potencia, lo cual es esencial para la dispensación efectiva de comida en el dispensador y para evitar problemas de desempeño.

    \item Se sugiere adaptar el diseño del dispensador en función del tipo de animal al que va dirigido. Ajustar el tamaño de los recipientes y el mecanismo de dispensación según las necesidades específicas del animal puede mejorar la eficacia del sistema y proporcionar una solución más adecuada para cada tipo de mascota.
\end{itemize} 